\documentclass[11pt]{article}

\usepackage[margin=0.85in]{geometry}
\usepackage[hidelinks]{hyperref}
\usepackage{array}
\usepackage{longtable}
\usepackage{enumitem}
\usepackage{titlesec}

\setlength{\parindent}{0pt}
\setlength{\parskip}{0.55em}
\emergencystretch=2em
\setlist[itemize]{itemsep=0.15em, topsep=0.25em}
\setlist[enumerate]{itemsep=0.15em, topsep=0.25em}
\titleformat{\section}{\large\bfseries}{\thesection}{0.6em}{}
\titleformat{\subsection}{\normalsize\bfseries}{\thesubsection}{0.6em}{}

\begin{document}

\begin{center}
    {\LARGE \textbf{Project Report}}\\
    \vspace{0.3em}
    {\large Comparison between PostgreSQL and Neo4j for AI Research Paper Retrieval and Knowledge Graph Analysis}\\
    \vspace{0.6em}
    Data Management 2025/2026\\
    Last updated: May 14, 2026
\end{center}

\section{Group Members}

\begin{itemize}
    \item Hamza Abdul Kader, matricola 2230150
    \item Leonardo Marasca, matricola 2217140
\end{itemize}

\section{Project Goal}

The goal of this project is to compare a relational DBMS and a graph database using the same AI research paper dataset.
We model and analyze OpenAlex research paper data in PostgreSQL and Neo4j.

The comparison focuses on:

\begin{itemize}
    \item data modeling
    \item query complexity
    \item expressiveness
    \item execution time
    \item suitability for relationship-heavy analysis
\end{itemize}

The project is related to Artificial Intelligence and Retrieval-Augmented Generation because the selected dataset focuses on AI, Machine Learning, Large Language Models, and RAG-related research papers.

\section{Dataset}

The dataset source is OpenAlex.

\begin{itemize}
    \item Main website: \url{https://openalex.org/}
    \item Dataset documentation: \href{https://help.openalex.org/hc/en-us/articles/24397285563671-About-the-data}{OpenAlex About the data}
\end{itemize}

We access the data through the OpenAlex API and collect a selected subset of research papers related to:

\begin{itemize}
    \item retrieval augmented generation
    \item large language models
    \item artificial intelligence
    \item machine learning
\end{itemize}

We do not use the full OpenAlex dataset because it is very large. Instead, we collect a controlled subset that is large enough for meaningful analysis but small enough to manage in a student project.

\section{Current Dataset Sample}

The current sample was collected with:

\begin{verbatim}
python src/collect_openalex.py --per-term 75 --per-page 75
\end{verbatim}

\begin{center}
\begin{tabular}{|l|r|}
\hline
\textbf{Entity} & \textbf{Count} \\
\hline
Papers & 299 \\
Authors & 1966 \\
Topics & 1104 \\
Paper-author relationships & 2077 \\
Paper-topic relationships & 4288 \\
Citation relationships & 392 \\
\hline
\end{tabular}
\end{center}

Generated data files are stored locally under \texttt{data/raw/} and \texttt{data/processed/}. These generated files are ignored by git because they can be regenerated from the OpenAlex API.

\section{Repository Structure}

\begin{verbatim}
src/
  collect_openalex.py        OpenAlex data collection and normalization
  benchmark_queries.py       Query timing benchmark runner
  generate_benchmark_charts.py

data/raw/                    Raw OpenAlex API output, ignored by git
data/processed/              Generated CSV files, ignored by git

database/postgres/           PostgreSQL schema, loading script, SQL queries
database/neo4j/              Neo4j constraints, import script, Cypher queries

benchmarks/
  queries.json               Paired SQL/Cypher benchmark queries
  results/benchmark_results.csv

docs/
  project_proposal.pdf
  project_roadmap.pdf
  data_model.md
  query_results.md
  benchmark_results.md
  project_report.md
  project_report.pdf
\end{verbatim}

\section{Data Collection Pipeline}

The script \texttt{src/collect\_openalex.py} collects data from the OpenAlex Works API.

Main steps:

\begin{enumerate}
    \item Search OpenAlex works using AI-related search terms.
    \item Use cursor pagination to retrieve results.
    \item Store raw API results as JSONL.
    \item Normalize nested OpenAlex data into CSV files.
    \item Reconstruct abstracts from OpenAlex abstract inverted indexes when available.
    \item Keep only citation links where both papers are included in the selected subset.
    \item Remove duplicate paper-author relationships before database loading.
\end{enumerate}

The output is designed so the exact same CSV files can be loaded into both PostgreSQL and Neo4j.

\section{PostgreSQL Model}

PostgreSQL uses a normalized relational model.

\begin{itemize}
    \item \texttt{papers}
    \item \texttt{authors}
    \item \texttt{topics}
    \item \texttt{paper\_authors}
    \item \texttt{paper\_topics}
    \item \texttt{citations}
\end{itemize}

Papers, authors, and topics are stored as entity tables. Many-to-many relationships are represented with bridge tables. Citations are represented as pairs of paper IDs. Indexes are created for common query patterns.

This model is strong for structured aggregation, filtering, sorting, and traditional SQL analysis.

\section{Neo4j Model}

Neo4j uses a graph model.

\begin{itemize}
    \item Nodes: \texttt{Paper}, \texttt{Author}, \texttt{Topic}
    \item Relationships: \texttt{AUTHORED}, \texttt{HAS\_TOPIC}, \texttt{CITES}
\end{itemize}

Example graph patterns:

\begin{verbatim}
(Author)-[:AUTHORED]->(Paper)
(Paper)-[:HAS_TOPIC]->(Topic)
(Paper)-[:CITES]->(Paper)
\end{verbatim}

This model is natural for relationship-heavy questions, such as author collaboration, citation paths, shared topics, and author citation networks.

\section{Docker Setup}

The project uses Docker Compose to run both databases locally:

\begin{itemize}
    \item PostgreSQL 16
    \item Neo4j 5
\end{itemize}

Start the databases:

\begin{verbatim}
docker compose up -d
\end{verbatim}

Stop the databases:

\begin{verbatim}
docker compose down
\end{verbatim}

Neo4j browser:

\begin{verbatim}
http://localhost:7474
username: neo4j
password: openalex123
\end{verbatim}

\section{Database Loading}

PostgreSQL loading:

\begin{verbatim}
docker exec openalex-postgres psql -U openalex -d openalex_ai -f /sql/schema.sql
docker exec openalex-postgres psql -U openalex -d openalex_ai -f /sql/load.sql
\end{verbatim}

Neo4j loading:

\begin{verbatim}
docker exec openalex-neo4j cypher-shell -u neo4j -p openalex123 \
  -f /cypher/constraints.cypher
docker exec openalex-neo4j cypher-shell -u neo4j -p openalex123 \
  -f /cypher/import.cypher
\end{verbatim}

Both databases were loaded successfully with the same generated CSV files.

\section{Validation}

We verified that PostgreSQL and Neo4j contain matching data counts.

\begin{center}
\begin{tabular}{|l|r|r|}
\hline
\textbf{Data element} & \textbf{PostgreSQL} & \textbf{Neo4j} \\
\hline
Papers & 299 & 299 \\
Authors & 1966 & 1966 \\
Topics & 1104 & 1104 \\
Paper-author / AUTHORED & 2077 & 2077 \\
Paper-topic / HAS\_TOPIC & 4288 & 4288 \\
Citations / CITES & 392 & 392 \\
\hline
\end{tabular}
\end{center}

This confirms that the same dataset was loaded consistently into both database systems.

\section{Comparison Queries}

We created paired SQL and Cypher queries for the following analyses:

\begin{enumerate}
    \item Most cited papers
    \item Authors with the most papers
    \item Most frequent topics
    \item Author collaboration pairs
    \item Citation links inside the selected subset
    \item RAG-related papers by title or abstract
    \item Two-hop citation paths
    \item Authors connected through shared topics
    \item Papers sharing cited references
    \item Citation paths up to three hops
    \item Author citation network
\end{enumerate}

The initial query results are documented in \texttt{docs/query\_results.md}. The benchmarked queries are stored in \texttt{benchmarks/queries.json}.

\section{Benchmark Methodology}

The benchmark script is:

\begin{verbatim}
python src/benchmark_queries.py --runs 5 --warmups 1
\end{verbatim}

Benchmark methodology:

\begin{itemize}
    \item one warmup run per query/system pair
    \item five measured runs per query/system pair
    \item PostgreSQL timing uses \texttt{EXPLAIN (ANALYZE, FORMAT JSON)}
    \item Neo4j timing uses \texttt{PROFILE} through \texttt{cypher-shell --format verbose}
    \item raw benchmark results are saved in \texttt{benchmarks/results/benchmark\_results.csv}
    \item benchmark summary is saved in \texttt{docs/benchmark\_results.md}
\end{itemize}

\section{Benchmark Results}

\small
\setlength{\tabcolsep}{3pt}
\begin{longtable}{|p{0.39\textwidth}|r|r|p{0.15\textwidth}|}
\hline
\textbf{Query} & \textbf{PostgreSQL avg ms} & \textbf{Neo4j avg ms} & \textbf{Faster system} \\
\hline
\endhead
Most cited papers & 0.058 & 4.200 & PostgreSQL \\
\hline
Authors with the most papers & 2.422 & 8.600 & PostgreSQL \\
\hline
Most frequent topics & 4.762 & 15.200 & PostgreSQL \\
\hline
Author collaboration pairs & 44.358 & 65.000 & PostgreSQL \\
\hline
Citation links inside the subset & 0.183 & 4.400 & PostgreSQL \\
\hline
RAG-related papers & 3.515 & 8.600 & PostgreSQL \\
\hline
Two-hop citation paths & 1.577 & 8.400 & PostgreSQL \\
\hline
Authors connected through shared topics & 309.447 & 586.600 & PostgreSQL \\
\hline
Papers sharing cited references & 2.171 & 7.800 & PostgreSQL \\
\hline
Citation paths up to three hops & 5.104 & 12.400 & PostgreSQL \\
\hline
Author citation network & 48.214 & 37.400 & Neo4j \\
\hline
\end{longtable}
\normalsize
\setlength{\tabcolsep}{6pt}

At the current dataset size, PostgreSQL is faster on most measured queries, especially simple ranking and aggregation queries. Neo4j is faster on the author citation network query, where the query follows a natural author-to-paper-to-paper-to-author graph pattern.

This does not mean one system is always better than the other. The dataset is still small, and performance depends on indexes, cache state, query shape, and graph density.

For this project, the strongest conclusion is that PostgreSQL is very efficient for relational aggregation, while Neo4j is often more readable and conceptually direct for graph-shaped analysis.

\section{Benchmark Charts}

The project also contains SVG charts generated from the benchmark CSV:

\begin{itemize}
    \item \texttt{docs/figures/benchmark\_average\_times.svg}
    \item \texttt{docs/figures/benchmark\_relative\_time.svg}
\end{itemize}

They can be regenerated with:

\begin{verbatim}
python src/generate_benchmark_charts.py
\end{verbatim}

These SVG files are useful for the final slides and for visual inspection on GitHub.

\section{Technical Issue Fixed}

During PostgreSQL loading, we found that OpenAlex can contain duplicate author entries for the same paper. This caused a duplicate primary key error in \texttt{paper\_authors}.

Fix:

\begin{itemize}
    \item changed the normalization logic in \texttt{src/collect\_openalex.py}
    \item de-duplicated paper-author relationships by \texttt{(paper\_id, author\_id)}
    \item regenerated CSV files
    \item reloaded PostgreSQL successfully
\end{itemize}

This is useful to mention in the final presentation because it shows a realistic data cleaning issue.

\section{Current Git History}

\begin{verbatim}
7bb7f82 Add graph-specific benchmark queries
6bc2e6b Add benchmark visualization charts
19d3df9 Add query benchmarking results
4626f8f Add living project report
e4a7b0a Load databases and record first query results
b30eb6b Add OpenAlex data pipeline scaffold
a465622 Initialize project documentation
\end{verbatim}

\section{Current Status}

Completed:

\begin{itemize}
    \item project proposal approved
    \item repository created and pushed to GitHub
    \item project roadmap written
    \item OpenAlex data collector implemented
    \item PostgreSQL schema created
    \item Neo4j import model created
    \item Docker Compose setup created
    \item first real OpenAlex sample collected
    \item PostgreSQL loaded
    \item Neo4j loaded
    \item equivalent SQL/Cypher queries executed
    \item benchmark script added
    \item timing benchmark executed
    \item benchmark results documented
    \item benchmark charts generated
    \item graph-specific query pairs added
\end{itemize}

Next:

\begin{itemize}
    \item optionally increase the dataset size and rerun benchmarks
    \item start shaping the final presentation structure
    \item write the presentation storyline and slide outline
\end{itemize}

\section{Update Log}

\subsection*{May 14, 2026}

\begin{itemize}
    \item Added four graph-specific query pairs: author shared topics, papers sharing cited references, citation paths up to three hops, and author citation network.
    \item Updated PostgreSQL and Neo4j query files.
    \item Reran the benchmark with 11 query pairs.
    \item Regenerated benchmark charts.
    \item Updated benchmark results and report interpretation.
    \item Added presentation-ready benchmark charts.
    \item Added repeatable benchmark infrastructure.
    \item Created the detailed living project report.
    \item Loaded PostgreSQL and Neo4j with the same OpenAlex CSV files.
    \item Verified matching counts across both systems.
    \item Fixed duplicate paper-author relationships in the collector.
\end{itemize}

\subsection*{May 12, 2026}

\begin{itemize}
    \item Project proposal was approved by Professor Roberto Maria Delfino.
\end{itemize}

\subsection*{May 11, 2026}

\begin{itemize}
    \item Created the project proposal PDF.
    \item Created the roadmap PDF.
    \item Created the GitHub repository documentation.
\end{itemize}

\end{document}
